\begin{recipe}
[% 
    preparationtime = {\unit[15]{min}},
    bakingtime = {\unit[55]{min}},
    bakingtemperature = {\protect\bakingtemperature{
        topbottomheat=\unit[180]{\textcelcius},
        fanoven=\unit[160]{\textcelcius}
    }},
    portion = {\portion{10}},
    source = {Eigene Kreation aus mehreren Online Rezepten},
    calory = {245kcal | 4g Protein | 33g KH | 10g Fett},
    price = {0,48€},
    created = {10.09.2025},
    rating = {4},
]
{Bananenbrot}

%\graph{% pictures
%    big=pic/dessert/09_klassisches-bananenbrot
%}

\introduction{%
    Wenn man mal wieder braune Bananen übrig hat, die man nicht mehr essen will\ldots
}

\ingredients{
    3(--4) & Reife Bananen\index{Obst!Bananen}\\
    80 ml & Neutrales Öl\index{Öle, Saucen \& Würzmittel!Pflanzenöl}\\
    110 g & Brauner Zucker\index{Backzutaten \& Süßes!Zucker}\\
    2 & Eierersatz\index{Eier!Eierersatz}\\
    200 g & Weizenmehl\index{Getreide, Pasta \& Brot!Weizenmehl}\\
    60 g & Schokoladendrops\index{Backzutaten \& Süßes!Schokolade}\\
    3 TL & Backpulver\index{Backzutaten \& Süßes!Backpulver}\\
    1 Prise & Salz\index{Kräuter \& Gewürze!Salz}\\
    1 Pck & Vanillepulver (vegan)\index{Backzutaten \& Süßes!Vanille}\\
    1 Prise & Zimt\index{Kräuter \& Gewürze!Zimt}
}

\preparation{%
    \step%
    Backofen auf 160C Umluft vorheizen. Eine Kastenform einfetten oder mit Backpapier auslegen.
    
    \step%
    3 Bananen mit einer Gabel fein zerdrücken.\\

    \step%
    Eierersatz nach Angabe anmischen, mit Bananen, Öl und Zucker vermischen,.
        
    \step%
    Mehl, Backpulver, Salz, Schokolade und Vanille mischen und unter die Eiermasse rühren.
    
    \step%
    Etwa 55 Minuten backen. Vor dem Anschneiden vollständig auskühlen lassen.
}

\suggestion[Deko]{%
    Hat man eine vierte Banane, halbiert man diese seitlich und legt sie auf die Teigmasse mit Schnittfläche nach oben.
}

\hint{
    Die Online Rezepte mit Apfelmus, Kichererbsenwasser etc funktionieren auch, aber aus eigenen Erfahrung finde ich Eierersatzpulver (vom DM) am einfachsten und besten.
}

\end{recipe}